\documentclass[10pt]{article}

% Specify the margins and text size.
\setlength{\textwidth}{6.4in}
\setlength{\textheight}{9.5in}
\setlength{\oddsidemargin}{0pt}
\setlength{\evensidemargin}{0pt}
\setlength{\topmargin}{0pt}
\setlength{\hoffset}{.05in}
\setlength{\voffset}{-1in}

\setlength{\parskip}{5pt}
\setlength{\parindent}{0pt}

% Load some fonts and symbol packages
\usepackage{latexsym}
\usepackage{pifont}       % contains 'star' symbol for counterinsurgency handbook title
\usepackage{yfonts} 
\usepackage{amsmath}
\usepackage{amsfonts}

\usepackage{graphicx}     % actually, this is loaded by pstricks
\usepackage[T1]{fontenc}
\usepackage{ifthen}
\usepackage{pstricks,pst-grad,pst-text,pst-node,multido,pst-plot,calc,pst-3dplot}
%\usepackage[all]{xy}
%\usepackage{animate}

% The hyperref package inserts links into the pdf.
\definecolor{MyLinkColor}{rgb}{.1,.2,1}
\definecolor{MyCiteColor}{rgb}{.1,1,.2}
\definecolor{MyURLColor}{rgb}{.4,.4,.4}
\usepackage[backref=true,pagebackref=false,hyperindex,colorlinks=true,
  linkcolor=MyLinkColor,urlcolor=MyURLColor]{hyperref}


% The tweaklist package is something I found on the web.  It provides a simple interface
% for making changes to spacing used in the itemize and enumerate environments.  Comment
% this out if you don't care to use tweaklist.
\usepackage{tweaklist}
\renewcommand{\itemhook}{\setlength{\parskip}{2pt}\setlength{\parsep}%
{1pt}\setlength{\topsep}{0pt}\setlength{\itemsep}{0pt}}

\newcommand{\U}{\underline{\hspace{5pt}}}

\usepackage{listings}
\newcommand{\Z}{\hphantom{0}}

\begin{document}
\pagestyle{empty}
\lstset{language=R, showspaces=false, showstringspaces=false}
\href{http://www.su.edu}{\includegraphics[height=1.75cm]{sulogo.eps}}
\vspace{-1.79cm}

{{\ }\hfill\small
\begin{tabular}{cl}
& Math 207\\
& Introduction to Statistics\\
\end{tabular}
}

\setlength{\baselineskip}{1.05\baselineskip}
\medskip

\begin{center}
\textbf{\large  Worksheet:  Tests of Significance}
\end{center}
\smallskip

\newcommand{\HH}{\hspace{20pt}\hphantom{a) } }

1. A die is rolled $n$ times and the total number of spots counted.
For each of the following cases, determine if the result can be explained as a chance
variation.  Calculate the test statistic $z$ and the $P$-value
(i.e.,~the chance of getting a test statistic 
  as extreme or more extreme than the one given).

\hspace{20pt} a) $n=100$, $\mbox{sum}=367$
\vspace{1in}

\hspace{20pt} b) $n=100$, $\mbox{sum}=316$
\vspace{1in}

\hspace{20pt} c) $n=100$, $\mbox{sum}=355$
\vspace{1in}

\hspace{20pt} d) $n=100$, $\mbox{sum}=384$
\vspace{1in}

\hspace{20pt} e) $n=100$, $\mbox{sum}=450$
\vspace{1in}

\hspace{20pt} f) $n=1,000$, $\mbox{sum}=3554$
\vspace{1in}

\hspace{20pt} g) $n=1,000$, $\mbox{sum}=3800$

\vfill
\eject
{\ }

2. A coin is flipped $n$ times and the total fraction of tails counted.
For each of the following cases, determine if the result can be explained as a chance
variation.  Calculate the test statistic $z$ and the $P$-value
(i.e.,~the chance of getting a test statistic 
  as extreme or more extreme than the one given).

\hspace{20pt} a) $n=100$, $\mbox{fraction of tails}=0.40$
\vspace{1.1in}

\hspace{20pt} b) $n=100$, $\mbox{fraction of tails}=0.45$
\vspace{1.1in}

\hspace{20pt} c) $n=100$, $\mbox{fraction of tails}=0.35$
\vspace{1.1in}

\hspace{20pt} d) $n=100$, $\mbox{fraction of tails}=0.55$
\vspace{1.1in}

\hspace{20pt} e) $n=100$, $\mbox{fraction of tails}=0.60$
\vspace{1.1in}

\hspace{20pt} f) $n=1,000$, $\mbox{fraction of tails}=0.45$
\vspace{1.1in}

\hspace{20pt} g) $n=1,000$, $\mbox{fraction of tails}=0.60$

\vfill
\eject

\end{document}


